\chapter{Clubhead Kinematics from Radar Velocities: A Screw-Theoretic How-To}
\label{app:screw}

This appendix develops, at implementation depth, the estimation layer that turns per-detection radar velocities into clubhead motion: what the current OpenFlight hardware can and cannot recover, why the twist (screw) representation of rigid-body motion is the natural formalism for the problem, and a step-by-step recipe for building the estimator on the IWR6843-class sensor described in \cref{app:hardware}.
It extends the EKF architecture of \cref{app:ekf} from point tracking (the ball) to rigid-body tracking (the club).

\section{Context: what the commercial systems actually recover}
\label{sec:screw-context}

TrackMan's club-data pipeline, as described in its documentation and patent family, tracks the clubhead ``from about knee height to impact,'' resolves the head's radar return into velocity components across multiple receivers, and reports the trajectory of the head's \emph{geometric center}~\cite{trackmanclubspeed,us10850179,trackmanoert}.
The physical basis is that a rotating, translating clubhead is not a point target: the toe moves up to $\sim$7\mph{} faster than the heel, so the club's return occupies a spread of Doppler bins, and multi-receiver phase interferometry (\cref{sec:interferometry}) locates each velocity component in angle.
What is fitted to those observations is a \emph{rigid-body motion model}, not an image; ``3D silhouette'' is marketing shorthand for that model fit.

\paragraph{Is screw theory used by the industry?}
There is no public evidence that TrackMan (or any launch-monitor vendor) formulates this fit in screw/twist coordinates: the patents describe point-trajectory tracking, Doppler-component decomposition, and silhouette correlation, without reference to the screw formalism~\cite{us8845442,us10850179}.
The instantaneous screw axis (ISA) \emph{is}, however, an established tool in golf biomechanics: Vena et al.\ applied ISA theory to optical motion capture of the swing in a two-part \emph{Sports Engineering} study, verifying that segment motion during the downswing is dominantly rotational about a well-defined moving axis (at least 71\% of marker velocity attributable to ISA rotation) and using ISA smoothness to characterize the kinematic sequence~\cite{vena2010a,vena2010b}.
The contribution proposed here --- estimating the club's twist \emph{directly from Doppler detections} rather than from marker positions --- is, to our knowledge, not published in the golf literature, although the underlying mathematics is standard in robotics~\cite{murrayliss} and radar micro-Doppler analysis~\cite{chenmicrodoppler}.
For OpenFlight this is an opportunity: the formalism is public-domain mathematics, distinct from the specific claimed pipelines in \cref{app:patents}.

\section{Screw theory in the minimum required dose}
\label{sec:screw-primer}

A rigid body's instantaneous motion is fully described by a \textbf{twist}
\begin{equation}
\label{eq:screw-twist}
\xi = (\vect{\omega},\, \vect{v}_O) \in \mathfrak{se}(3),
\end{equation}
where $\vect{\omega}$ is the angular velocity and $\vect{v}_O$ is the linear velocity of a chosen body reference point $O$.
The velocity of any body-fixed point at position $\vect{r}$ relative to $O$ is then
\begin{equation}
\label{eq:screw-pointvel}
\vect{v}(\vect{r}) = \vect{v}_O + \vect{\omega} \times \vect{r}.
\end{equation}
Chasles' theorem states that every such motion is instantaneously a rotation about, plus a translation along, a unique line in space: the \textbf{instantaneous screw axis} (ISA).
Its direction is $\hat{\vect{\omega}}$; a point on it is
\begin{equation}
\label{eq:screw-isa}
\vect{r}_{\mathrm{ISA}} = \frac{\vect{\omega} \times \vect{v}_O}{\lVert\vect{\omega}\rVert^{2}},
\qquad
h = \frac{\vect{\omega} \cdot \vect{v}_O}{\lVert\vect{\omega}\rVert^{2}},
\end{equation}
where the \textbf{pitch} $h$ is the translation per radian along the axis.
For a downswing near impact the club's motion is close to a pure rotation about a hub near the hands: the ISA passes near the grip and the pitch is small.
That geometric fact is both a physical insight (the ISA trajectory \emph{is} a rigorous definition of swing plane) and, later, a regularization prior.
Standard references: Murray, Li \& Sastry for the mathematics~\cite{murrayliss}; Vena et al.\ for golf-specific ISA practice and its error behavior~\cite{vena2010a}.

\section{The measurement model: Doppler is linear in the twist}
\label{sec:screw-measurement}

A radar detection assigns to some scattering center at known position $\vect{p}_i$ (from range and monopulse angle) a radial velocity $\dot d_i$ along the unit line of sight $\uvec{u}_i = \vect{p}_i / \lVert\vect{p}_i\rVert$.
Substituting \cref{eq:screw-pointvel} with $\vect{r}_i = \vect{p}_i - \vect{p}_O$:
\begin{equation}
\label{eq:screw-doppler}
\dot d_i
= \uvec{u}_i \cdot \vect{v}(\vect{r}_i)
= \underbrace{\uvec{u}_i}_{1\times3} \cdot\, \vect{v}_O
\;+\; \underbrace{(\vect{r}_i \times \uvec{u}_i)}_{1\times3} \cdot\, \vect{\omega}.
\end{equation}
This is the reciprocal product of the sight line (as a Pl\"ucker line) with the twist, and it is \textbf{exactly linear} in the six unknowns $(\vect{v}_O, \vect{\omega})$.
Each detection contributes one row of a linear system
\begin{equation}
\label{eq:screw-ls}
\vect{z} = A\,\xi + \vect{\epsilon},
\qquad
A_i = \bigl[\; \uvec{u}_i^{\mathsf T} \;\big|\; (\vect{r}_i \times \uvec{u}_i)^{\mathsf T} \;\bigr],
\end{equation}
so the per-frame twist estimate is weighted least squares --- no iteration, no linearization error, and a covariance $(A^{\mathsf T} W A)^{-1}\sigma^2$ for free.
This is the property that makes the screw formulation not merely elegant but \emph{operationally} correct for radar: the sensor's native observable is already a linear functional of the twist.
(Camera systems enter the same framework from the other side: a fiducial-tracked face gives an $SE(3)$ pose per frame, and the matrix logarithm of the frame-to-frame relative pose is a finite twist~\cite{murrayliss}; one downstream representation serves both modalities.)

\section{What the current hardware can observe}
\label{sec:screw-ceiling}

Apply \cref{eq:screw-ls} to each OpenFlight sensor to see the ceiling precisely.

\subsection{OPS243-A: one row, no position}
The OPS243-A has a single receive channel: it measures $\dot d_i$ but neither range nor angle, so $\vect{p}_i$ is unknown and \cref{eq:screw-ls} cannot be assembled.
What survives is the \emph{marginal distribution} of $\dot d$ over the head: at 30\,ksps with 128-sample segments the raw Doppler bin is $\approx$3.3\mph{} (\cref{app:hardware}), so the full toe--heel spread of a driver spans only $\sim$2 raw bins.
Software can and should still extract (\cref{app:radardsp}): the spread's midpoint or a fixed percentile as a stable club-speed reference (the peak is toe glint); the spread \emph{width} as a crude proxy for $\lVert\vect{\omega}\rVert$ projected on the line of sight; and head/shaft separation by gating the slow tail.
That is the honest ceiling of the current stack: a one-dimensional shadow of the twist, useful for de-biasing club speed, incapable of yielding path or attack angle.

\subsection{K-LD7: the right equation, the wrong operating envelope}
The K-LD7's dual receive patches provide per-bin monopulse angle --- rows of \cref{eq:screw-ls} in principle --- but its fastest setting has a 100\,km/h ($\approx$62\mph) unambiguous-velocity ceiling and $\sim$29\,ms frames (\cref{app:hardware}), so a driver head aliases and traverses more than a meter between frames.
It remains a ball-burst angle sensor, not a club tracker.

\subsection{IWR6843: all three ingredients in one package}
The IWR6843 supplies everything \cref{eq:screw-ls} needs, given custom chirp firmware (\cref{app:hardware}): range at 3.75\,cm resolution (separates head from shaft and arms), per-chirp Doppler with both span and resolution set by the chirp plan, and per-detection monopulse angle from 12 virtual antennas.
Representative chirp budget for the club problem:
\begin{center}\small
\begin{tabular}{@{}llll@{}}
\toprule
\textbf{Choice} & \textbf{Value} & \textbf{Consequence} & \textbf{Governing relation} \\
\midrule
Chirp repetition $T_c$ & 25\,\si{\micro\second} & $v_{\max} = \lambda/4T_c \approx 50$\,m/s (112\mph) & unambiguous velocity \\
Chirps per frame $N$ & 128 & $\Delta v = \lambda/2NT_c \approx 0.8$\,m/s & velocity resolution \\
Frame time & 3.2\,ms & $\sim$300\,Hz frame rate & swing-resolved motion \\
Bandwidth $B$ & 4\,GHz & $\Delta R = c/2B = 3.75$\,cm & 3--4 cells across the head \\
\bottomrule
\end{tabular}
\end{center}
At $\lambda = 5$\,mm the toe--heel spread ($\sim$3\,m/s) covers $\sim$4 velocity cells --- twice the OPS243-A's resolving power, with each cell additionally localized in range and angle.
The out-of-box SDK point cloud at 10--20\,Hz is useless here; the custom chirp configuration plus on-chip range-Doppler processing is the enabling (and largest) engineering task, as flagged in \cref{app:hardware}.

\section{Estimation recipe, step by step}
\label{sec:screw-recipe}

The following is the full pipeline, in execution order, for one swing.
Notation: frames indexed $k$, detections within a frame indexed $i$.

\subsection*{Step 1 --- Detect and localize}
Run 2D CFAR on each frame's range-Doppler map (OS-CFAR preferred over cell-averaging: club and ball are exactly the closely-spaced-targets case it is designed for, \cref{app:hardware}).
For each detection, compute monopulse azimuth/elevation from the virtual-array phase and form $\vect{p}_i = (R_i, \theta_i, \phi_i)$ in Cartesian sensor coordinates, with a per-detection covariance driven by SNR (angle variance $\propto 1/\mathrm{SNR}$; range variance from the window function).

\subsection*{Step 2 --- Segment the club}
Gate detections to the club by a joint predicate: range inside the swing corridor, radial speed in the club band (above waggle, below ball speeds), and --- once tracking has started --- Mahalanobis distance to the propagated track.
Reject the shaft by its distinct signature: shaft returns sit at lower speed and nearer range, forming an elongated cluster whose major axis points at the head cluster.
Body and arm clutter falls out on speed alone.

\subsection*{Step 3 --- Per-frame twist solve}
Choose the body reference point $O_k$ as the SNR-weighted centroid of the head cluster (this choice only re-parameterizes $\vect{v}_O$; the twist itself is frame-invariant).
Assemble \cref{eq:screw-ls} over the frame's $m_k$ detections and solve weighted least squares with a robust loss (Huber or Cauchy) on the residuals --- robustness is not optional, because \emph{glint migration} (specular points wandering across the curved metal head as aspect changes) violates the fixed-scatterer assumption at the few-centimeter level and manifests as heavy-tailed residuals.
Record $\hat\xi_k$ and its covariance $\Sigma_k = (A^{\mathsf T} W A)^{-1}$.

\subsection*{Step 4 --- Read the conditioning honestly}
Compute the SVD of the weighted $A$.
Expect, from a single vantage $\sim$2\,m away with detections spanning $\sim$10\,cm: one strong singular value (bulk radial velocity --- essentially club speed), one or two moderate ones (the transverse velocity component resolved by angle diversity, and the $\vect{\omega}$ combination that generates the Doppler spread), and the rest weak.
\textbf{Do not invert the weak directions per frame.}
Either truncate the SVD (solve only the well-conditioned subspace and carry the null-space explicitly) or defer those components entirely to the smoother of Step 5.
A per-frame 6-DOF solve that ``works'' numerically will hallucinate the unobservable components from noise; this is the single most likely implementation failure mode.

\subsection*{Step 5 --- Smooth the twist trajectory on $SE(3)$}
The rank deficiency resolves \emph{across} frames: as the head sweeps through the beam over $\sim$30\,ms ($\sim$10 frames at 300\,Hz), the sight-line geometry rotates and different frames constrain different twist combinations.
Estimate a smooth trajectory rather than independent snapshots:
state $x_k = (T_k \in SE(3),\, \xi_k)$, i.e., pose and twist;
process model $T_{k+1} = T_k \exp\!\bigl((\xi_k \Delta t)^{\wedge}\bigr)$ with a slowly varying twist (white-noise angular/linear acceleration, tuned to swing dynamics --- the head gains roughly 1\mph{} per millisecond late in the downswing);
measurements = the raw rows of \cref{eq:screw-doppler} (preferable to feeding the Step-3 point estimates, since rows carry their true information content);
run an extended/unscented filter with the standard Lie-group state handling~\cite{murrayliss}, then an RTS backward smoother, exactly as in \cref{app:ekf}.
Add two physically motivated priors, weighted weakly enough to be overruled by data:
\begin{enumerate}
\item \textbf{Low pitch}: penalize $h^2 = (\vect{\omega}\cdot\vect{v}_O)^2/\lVert\vect{\omega}\rVert^4$ --- the downswing is close to a pure rotation~\cite{vena2010a}.
\item \textbf{Hub proximity}: penalize distance of the ISA (\cref{eq:screw-isa}) from a broad prior region around the golfer's hands.
\end{enumerate}

\subsection*{Step 6 --- Evaluate at impact, then stop}
Rigidity holds before impact and fails during it ($\sim$500\,\si{\micro\second} of gross deformation), so fit only up to the last pre-impact frame and evaluate the smoothed twist at the impact timestamp (from the sound trigger, minus the acoustic delay budget of \cref{app:hardware}).
This mirrors the commercial convention: TrackMan defines club speed ``just prior to first contact''~\cite{trackmanclubspeed}.

\subsection*{Step 7 --- Project out the parameters}
All club-delivery parameters are now projections of one estimated object:
\begin{center}\small
\begin{tabular}{@{}p{3.6cm}p{9.6cm}@{}}
\toprule
\textbf{Parameter} & \textbf{Projection of the impact-time twist} \\
\midrule
Club speed & $\lVert \vect{v}_O + \vect{\omega}\times\vect{r}_{gc} \rVert$ at a \emph{declared} reference point $\vect{r}_{gc}$ (cluster centroid, or a per-club calibrated offset) \\
Club path & horizontal direction angle of that velocity vector vs.\ the target line \\
Attack angle & vertical direction angle of the same vector \\
Closure rate & component of $\vect{\omega}$ about the estimated shaft axis, in \si{\degree\per\second} --- the GCQuad-exclusive parameter (\cref{sec:camclub}), free here \\
Swing plane / direction & orientation of the ISA (\cref{eq:screw-isa}) and of the plane its trajectory sweeps \\
Low point & extrapolate the reference-point arc to its velocity-vertical zero \\
Gear-effect recoil & (post-impact, optional) the discontinuity in $\vect{\omega}$ across impact estimates the head's angular recoil, cross-checkable against \cref{eq:gear} \\
\bottomrule
\end{tabular}
\end{center}
Report each with the covariance propagated from the smoother, and tag provenance per \cref{tab:hierarchy}: with this pipeline, path and attack angle move from \emph{derived} to \emph{measured}, while face angle and impact location remain optics problems (\cref{sec:faceangle}).

\section{Validation plan}
\label{sec:screw-validation}

Validate in order of increasing realism, reusing the standards of \cref{ch:accuracy}:
\begin{enumerate}
\item \textbf{Synthetic}: simulate scatterers on a CAD clubhead following a recorded swing trajectory, generate detections with realistic SNR and glint migration, and verify the recovered twist against truth --- this exercises Steps 3--5 without hardware and quantifies the observability claims of Step 4.
\item \textbf{Pendulum rig}: a club swung as a physical pendulum has an analytically known, planar, fixed-axis twist (pitch exactly zero, ISA fixed); any bias in the pipeline shows up immediately.
\item \textbf{Cross-device}: compare club speed, path, and attack angle against the MLM2PRO per the protocol of \cref{ch:accuracy}, with the reference-point caveat of \cref{sec:radarclub} stated up front; sanity-gate every shot against the loft-dependent smash ceiling (\cref{sec:smash}).
\item \textbf{Internal consistency}: the measured path must satisfy the swing-geometry identity (path as a function of swing direction, swing plane, attack angle --- \cref{ch:impact}) within uncertainty; violations indicate reference-point drift or alignment error.
\end{enumerate}

\begin{implication}
The twist formulation costs nothing extra to adopt --- the per-frame solve is a small weighted least squares --- and buys three things the point-tracking alternative cannot: a principled, declared club-speed reference point (ending the largest inter-device disagreement), closure rate and a rigorous swing plane as free by-products, and honest covariances on every club parameter.
It should be specified as the estimation layer for the IWR6843 integration from day one, with the OPS243-A velocity-distribution analytics (\cref{app:radardsp}) as the degenerate 1D special case of the same model.
\end{implication}
